\begin{textsample}{POS Dim 3 – vem\_ed\_16 – Score 11.00 – vem\_ed\_16/t070.txt}  \label{ex:f3_pos_015}
VEm Virtual Exchange Medium Como vê as iniciativas de \textbf{cooperação} Sul-Sul nos Intercâmbios Virtuais?
É explícita, nas instituições latino-americanas, por exemplo, a resistência a \textbf{cooperação} com o Sul Global.
Uma \textbf{breve} \textbf{análise} do número de acordos de nossas IES com cada região do planeta poderá, facilmente, mostrar que os países do Norte, com suas instituições centenárias e suas ofertas tentadoras, estão na dianteira do interesse de nossos estudantes e pesquisadores como destino de intercâmbio.
A possibilidade dos Intercâmbios Virtuais entre instituições do Sul abre, pois, uma via a ser explorada para uma quebra inicial – que poderá ser seguida de um desejável equilíbrio – da primazia do Norte em nossas relações bilaterais.
As trocas virtuais permitirão o desenvolvimento de algo que é fundamental para a \textbf{cooperação} Sul-Sul, a saber, o conhecimento mútuo dos parceiros potenciais, de sua atuação na pesquisa, de sua cultura.
Somente esse conhecimento permitirá a construção de sólidas parcerias entre instituições do Sul, talvez com uma reciprocidade maior que aquela que temos observado nas relações com o Norte.
Comente a conquista do \textbf{selo} FAUBAI-BRaVE pelo CPS, em \textbf{outubro} de 2022.
O \textbf{selo} FAUBAI-BRaVE foi criado para atestar a participação das IES \textbf{membros} da FAUBAI nas nossas promoções de Intercâmbio Virtual nos moldes do Programa FAUBAI-BRaVE.
A atribuição do \textbf{selo} ao Centro Paula Souza, além de reconhecer a qualidade de seus projetos de Intercâmbio Virtual, representa também o reconhecimento da FAUBAI pela participação do CPS, juntamente com a UNESP e a UFPE, na experiência piloto que construiu esse programa.
Como as demais IES portadoras desse \textbf{selo}, o CPS poderá utilizá-lo para comprovar sua participação no programa diante dos parceiros que vier a \textbf{conquistar} e, também, terá direito a condições especiais nas promoções que o programa FAUBAI-BRaVE vier a lançar no futuro, como ofertas de parcerias estabelecidas pela FAUBAI, gratuidades ou descontos em formações específicas para a área, entre outros.
Como vislumbra o futuro da internacionalização?
O cenário macropolítico internacional aparenta um recrudescimento de ideias totalitárias e, por vezes, nacionalistas, contra o ideal democrático que faz transpor fronteiras.
Essa \textbf{impressão} acentuou-se no período da pandemia, quando as trocas, num primeiro momento, pareciam ser impossíveis.
Entretanto, podemos dizer sem exagero que o braço tecnológico da globalização se estabeleceu solidamente em todas as culturas, permitindo-nos, inclusive, enfrentar com menos sofrimento as agruras das quarentenas.
Espantosamente, as tecnologias que, pelo alto \textbf{custo} do acesso à Web e a equipamentos, sempre foram excludentes, tornaram-se um fator de \textbf{inclusão} de primeira ordem no mundo da educação, para nos restringirmos ao nosso campo.
Isso mostra que detemos um conhecimento capaz de nos fazer superar qualquer obstáculo, incluindo o negacionismo e o obscurantismo que compõem as receitas totalitárias.
Vejo, pois, o futuro da educação superior internacionalizada sempre mais adepto das trocas virtuais, que já provaram ser de grande utilidade, mesmo diante do mau uso que pode, por vezes, permeá-la.
Ter o mundo em nossa casa, com o recurso quase mágico das novas tecnologias, é uma oportunidade que não pode ser ignorada, pois, além dos contatos rápidos e constantes, oferece também uma via mais segura e econômica dada a necessidade de redução da emissão de carbono.
Não se \textbf{trata} de imaginar o futuro da internacionalização apenas em um metaverso crescente, mas sim de integrar essas possibilidades virtuais ao processo educacional e à produção científica.

% matched lemmas: análise, breve, conquistar, cooperação, custo, impressão, inclusão, membro, outubro, selo, tratar
\end{textsample}
